\documentclass[11pt]{article}

\usepackage[T1]{fontenc}
\usepackage[utf8]{inputenc}
\usepackage{lmodern}
\usepackage[margin=1in]{geometry}
\usepackage{amsmath,amssymb,amsthm,mathtools}
\usepackage{enumitem}
\usepackage{booktabs}
\usepackage{array}
\usepackage{longtable}
\usepackage{xcolor}
\usepackage{microtype}
\usepackage{etoolbox}
\usepackage{url}
\usepackage[colorlinks=true,linkcolor=blue!50!black,citecolor=green!35!black,urlcolor=blue!60!black]{hyperref}
\usepackage[capitalize,nameinlink]{cleveref}

% ---------------------------------------------------------------------------
% Theorem environments
% ---------------------------------------------------------------------------
\theoremstyle{plain}
\newtheorem{theorem}{Theorem}[section]
\newtheorem{lemma}[theorem]{Lemma}
\newtheorem{corollary}[theorem]{Corollary}
\newtheorem{proposition}[theorem]{Proposition}
\newtheorem{conjecture}[theorem]{Conjecture}
\newtheorem{question}[theorem]{Question}
\newtheorem{problem}[theorem]{Problem}
\theoremstyle{definition}
\newtheorem{definition}[theorem]{Definition}
\newtheorem{remark}[theorem]{Remark}
\newtheorem{example}[theorem]{Example}

% The Lean proof is the proof; the text gives a faithful outline of it.
\newenvironment{outline}[1][Outline of the formal proof]{\begin{proof}[#1]}{\end{proof}}
% Results that are not formalized get complete proofs, labelled as such.
\newenvironment{informalproof}[1][Proof (not formalized)]{\begin{proof}[#1]}{\end{proof}}

% ---------------------------------------------------------------------------
% Lean references
%   \lean{name}            a declaration
%   \formal{file}{names}   the formal counterpart of a numbered result: a file and a comma-separated list of declarations
%   \partformal{file}{names}{gap}
%                          a numbered result whose formal counterpart covers all but the
%                          part described in {gap}; the gap is proved or cited in the text
%   \notformalized         a numbered result that has no formal counterpart
%   \cited{where}          a result of an earlier paper used without proof: where it is proved,
%                          and whether it is Lean-verified there
% ---------------------------------------------------------------------------
% Lean names are typeset like paths, with line breaks allowed after "_", "." and "/".
\DeclareUrlCommand\leanname{\def\UrlFont{\ttfamily\upshape}}
\renewcommand{\UrlBreaks}{\do\_\do\.\do\/\do\-}
\newcommand{\lean}[1]{\leanname{#1}}
\newcommand{\leanfile}[1]{\leanname{#1}}
\newcommand{\leanlist}[1]{\begingroup\def\leansep{}%
  \renewcommand*{\do}[1]{\leansep\leanname{##1}\def\leansep{, }}\docsvlist{#1}\endgroup}
\newcommand{\formal}[2]{\par\noindent{\small\upshape\raggedright\textsf{Lean:} \leanlist{#2}\unskip
  \nobreak\hfil\penalty50\hskip1em\hbox{}\nobreak\hfill[\leanfile{#1}]\parfillskip=0pt\par}%
  \smallskip\noindent\ignorespaces}
\newcommand{\partformal}[3]{\par\noindent{\small\upshape\raggedright\textsf{Lean (partial):} \leanlist{#2}\unskip
  \nobreak\hfil\penalty50\hskip1em\hbox{}\nobreak\hfill[\leanfile{#1}]\parfillskip=0pt\par}%
  \par\noindent{\small\upshape\textsf{Not formalized:} #3}\par\smallskip\noindent\ignorespaces}
\newcommand{\notformalized}{\par\noindent{\small\upshape\textsf{Not formalized;} a complete proof is given
  below.}\par\smallskip\noindent\ignorespaces}
\newcommand{\cited}[1]{\par\noindent{\small\upshape\textsf{Cited:} #1}\par\smallskip\noindent\ignorespaces}

% ---------------------------------------------------------------------------
% Notation
% ---------------------------------------------------------------------------
\newcommand{\N}{\mathbb{N}}
\newcommand{\Z}{\mathbb{Z}}
\newcommand{\Q}{\mathbb{Q}}
\newcommand{\R}{\mathbb{R}}
\newcommand{\C}{\mathbb{C}}
\newcommand{\F}{\mathbb{F}}
\newcommand{\T}{\mathsf{T}}
\newcommand{\clog}[1]{\lceil \log_2 #1 \rceil}
\newcommand{\abs}[1]{\lvert #1 \rvert}
\newcommand{\card}[1]{\lvert #1 \rvert}
\newcommand{\eps}{\varepsilon}
\newcommand{\prev}{\mathrm{prev}}
\newcommand{\cost}{\mathrm{cost}}
\newcommand{\gn}{\mathrm{gn}}
\newcommand{\Sh}{\mathrm{Sh}}
\newcommand{\cl}{\mathrm{cl}}
\newcommand{\Aut}{\mathrm{Aut}}
\newcommand{\minrk}{\mathrm{minrk}}
\newcommand{\MAIS}{\mathrm{MAIS}}
\newcommand{\Crux}{\mathrm{Crux}}
\newcommand{\TCP}{\mathrm{TCP}}
\DeclareMathOperator{\rank}{rank}
\DeclareMathOperator{\dist}{dist}
\DeclareMathOperator{\poly}{poly}


\title{Tape Transposition and the Matrix Coding Conjecture}
\author{Hans Brende \and Claude Fable 5.2 (Anthropic) \and GPT-6 Astra (OpenAI)}
\date{November 2026}

\begin{document}

\maketitle

\begin{abstract}
This paper is the companion of the first paper of the series, \emph{Matrix Transposition and the Tape Coding
Conjecture}. That paper reduced the question whether a multitape Turing machine can transpose a $d\times d$ matrix of
$\lceil\log_2d\rceil$-bit entries in linear time to a conjecture about network codes on the time--space diagrams of
tapes. Here the implication runs the other way: a conjecture about arrays alone, with no machine, no time and no
network in its statement, implies the Tape Coding Conjecture.

A \emph{whole-row code} of window~$w$ is a set of $d\times d$ arrays in which every column is a function of at most
$w$ rows. The symmetric arrays are such a code, with $w=1$ and about half of all entries free. The \emph{Matrix Coding
Conjecture} says that there is a $\delta>0$ such that, for every fixed~$w$ and all large~$d$, a whole-row code of
window~$w$ has at most $2^{(1-\delta)L}$ elements, where $L$ is the number of bits of an array. We prove that the
Matrix Coding Conjecture implies the Tape Coding Conjecture in its robust form, and therefore that no multitape Turing
machine, oblivious or not, would transpose in linear time. A weak form, in which the deficit $\delta$ may shrink slowly
as the window grows and only one size $d$ beyond every threshold is needed, already suffices. For oblivious machines the
reduction is direct: after a cheap cut of the run, the bits on which the outputs of almost every row of the transpose
depend lie in boundedly many rows of the input.

The conjecture is open, and the rest of the paper is about it. Every whole-row code contains a large code in which rows
are also functions of columns, on a symmetric relation. Every bound obtained from an acyclic set of entries is the
bound of an order of the rows, and such bounds, with or without transported copies of the recovery rules, cannot give
a deficit that is independent of the window. Entropy inequalities can go further: on the incidence relation of the
Fano plane the best acyclic bound is $41$ entries of $49$ and the Shannon value is $119/3$, and the gap persists, at more
than one percent of all entries, on every lift whose helper rows respect the incidences of the plane and which has a
mixing property. Codes that are invariant under the shifts of a
pairing of the entries, with $b$ pairs omitted, have rate at most $\tfrac12+b/d^2+o(1)$ over every finite abelian group
and for every relation of window~$w$, with nonlinear recovery allowed.

The theorems above are formalized in Lean~4, with the standard axioms only, against the frozen statement of the first
paper. Proved in the text and not formalized are two existence statements (relations of bounded window with the mixing
property, and mixing lifts of the Fano relation), the code that attains the Fano optimum, and one further result: the
exclusion of oblivious runs whose time diagram a cheap cut splits into pieces with few outputs.

\end{abstract}

\clearpage

% Contributions checked against the research record and confirmed by the authors. The clauses on
% \cref{prop:matchings} and on the review of the final text were added afterwards from the commit history.
\begin{quote}\small
\textbf{Contributions.} Hans Brende chose the problem, fixed the verification standard and directed the research. He
asked how the question about arrays relates to the Tape Coding Conjecture of the first paper, which led to
\cref{thm:T}, and asked for the weakest sufficient form, which led to \cref{thm:F}. Claude Fable~5.2, a large language
model developed by Anthropic, proposed the row-spread statement of \cref{sec:rows} with a first proof, and with it the
question about whole-row codes as the single target of the programme; proposed the two-sided form of
\cref{sec:normal}, the route through block cuts of \cref{sec:tape}, the grid relation of \cref{sec:entropy} and the
local-view lemma behind \cref{prop:hyperfinite}; proved \cref{prop:matchings} in its first form; reviewed the
statements against their formal counterparts; and wrote this paper apart from its proofs. GPT-6 Astra, a large
language model developed by OpenAI, proved \cref{thm:W,thm:T,thm:F}
in their final form; found and proved the translated intersection of \cref{sec:normal}, the barriers of
\cref{sec:order}, the entropy bounds, the three-lane code and the lift theorem of \cref{sec:entropy}, and the
invariance theorems of \cref{sec:invariance}; extended \cref{prop:matchings} to windows that grow with~$d$; found
the repairs that several of the proposed arguments needed; wrote the Lean formalization; wrote the proofs, the
Lean index and the verification appendix of this paper; and reviewed its final text.
\end{quote}

\tableofcontents
\clearpage

\section{Introduction}\label{sec:intro}

\subsection{The problem}\label{sec:intro-problem}

A $d\times d$ matrix with entries of $b(d)=\max(1,\lceil\log_2d\rceil)$ bits is written on the input tape of a
deterministic multitape Turing machine, row after row, as a word of $L=d^2b(d)$ bits. The machine has to write the
transpose. Statement~(T) says that some machine does this within $c(L+1)$ steps for a constant~$c$. The best known
machines need of the order of $L\log d$ steps, apart from a small saving in the exponent of the logarithm claimed in an
unrefereed preprint~\cite{OpenAI2026mult}, and no lower bound above linear is known for any explicit problem on machines with two work tapes. A
linear-time multiplier of integers would give a linear-time transposer~\cite{HvdH2025}, so a proof that (T) fails would
also bound multiplication from below. The formal statement is the Lean proposition
\lean{Transpose.LinearTimeTransposable} of the first paper~\cite{FirstPaper2026}, which is frozen: this paper proves
theorems about it and does not change it.

\subsection{Two conjectures}\label{sec:intro-two}

The first paper replaced (T) by a statement about codes. The run of a machine is drawn as a time--space diagram: the
positions of the heads at every step. Cutting the tapes into blocks gives a network whose nodes are the visits of a
head to a block, and a linear-time transposer yields a network code on it that delivers a piece of every row block to
every column block. The \emph{Tape Coding Conjecture} says that no such code exists on the block network of any
diagram of linear length; it implies that (T) fails~\cite[Conjecture~9.1 and Theorem~9.2]{FirstPaper2026}. Its
statement has no transition table in it, but it keeps the diagram: time, the shape of tapes, and the capacity of
every link.

The conjecture of this paper keeps none of these. Let $\rho$ assign to every column index a set of at most $w$ row
indices. A set $C$ of $d\times d$ arrays is a \emph{whole-row code for~$\rho$} if two arrays of $C$ that agree on the
rows $\rho(j)$ agree on column~$j$, for every~$j$ (\cref{def:wholerow}). For $\rho(j)=\{j\}$ the symmetric arrays are
such a code: column $j$ is row~$j$. They have $d(d+1)/2$ free entries, a little more than half. We call the fraction
of free entries the \emph{rate}.

\begin{conjecture}[Matrix Coding Conjecture]\label{conj:MCC}
There is a $\delta>0$ such that for every $w\ge1$ there is a $d_0$ with the following property. For every $d\ge d_0$
and every $\rho$ of window at most~$w$, every whole-row code for~$\rho$ with entries of $b(d)$ bits has at most
$2^{(1-\delta)L}$ elements.
\end{conjecture}

The conjecture asks only that the rate stay away from one. We expect much more, namely that the rate tends to one half
at every fixed window (the \emph{half-rate conjecture}, \cref{conj:half}), which is the form in which the synthesis
paper of the series asked the question for a fixed alphabet~\cite[Problem~9.4(c)]{SeventeenthPaper2026}. What is known
is the bound $w/(w+1)$ for the rate, and one half for $w\le2$~\cite{SeventeenthPaper2026}.

\subsection{Results}\label{sec:intro-results}

The first three theorems are the reduction. The other four are about the conjecture.

\begin{itemize}[leftmargin=1.5em]
\item \textbf{Theorem~W (whole rows; \cref{thm:W}).} In any run of linear length, a cut of bounded depth can be
extended, by revealing at most $\eps L$ further bits, so that the input bits on which the outputs of all but a small
fraction of the rows of the transpose depend lie in boundedly many rows of the input, the bound depending on $\eps$
and on the resources of the machine and not on~$d$. The theorem is about the motion of the heads alone: it assumes
neither that the machine is oblivious nor that it is correct. For oblivious machines, whose runs have cheap cuts of
bounded depth, it follows that the Matrix Coding Conjecture excludes linear-time transposers (\cref{cor:W-obl}). This is what turns a statement about runs into a statement about arrays.
\item \textbf{Theorem~T (tape coding; \cref{thm:T}).} The Matrix Coding Conjecture implies the robust form of the Tape
Coding Conjecture. With the reduction of the first paper it follows that the Matrix Coding Conjecture implies that
(T) fails, for all machines.
\item \textbf{Theorem~F (the weak form; \cref{thm:F}).} A weak form of the conjecture (\cref{conj:weak}) already
implies that (T) fails. It asks for a bound at one size $d$ beyond every threshold, and lets the deficit decrease
with the window; \cref{rem:weak-demand} says how fast it may decrease.
\item \textbf{Theorem~N (normal form; \cref{thm:N}).} Every whole-row code $C$ contains, after a translation, a code
with at least $\card C^2/2^{L}$ elements in which rows are functions of columns as well, and the relation may be taken
symmetric and reflexive at a small cost. So the conjecture may be attacked for such codes only.
\item \textbf{Theorem~O (order; \cref{thm:O}).} The largest acyclic set of entries has exactly $d^2$ minus the least
cost of an order of the rows, and every acyclic set lies in the one given by some order. Recovery rules of
transported copies of a code do not improve this bound, and on relations with a mixing property it leaves a deficit
that tends to zero as the window grows (such relations exist by \cref{prop:mixing}, which is not formalized). So no argument of this kind proves the conjecture.
\item \textbf{Theorem~E (entropy; \cref{thm:E}).} On the incidence relation of the Fano plane every polymatroid that
obeys the recovery constraints has value at most $119/3$, a code attains this, and the best acyclic bound is~$41$. A
$9\times6$ grid relation without mutual triangles has value at most $204/5$ against~$41$. Every lift of the Fano
relation whose helper rows respect the incidences and which mixes at one percent has a gap of more than $d^2/100$
between the two bounds. Such lifts exist with both degrees at most $3600$ and girth tending to infinity
(\cref{prop:lifts}, not formalized). So bounds beyond order exist at a fixed window; these relations already have a
constant deficit from order.
\item \textbf{Theorem~I (invariance; \cref{thm:I}).} A whole-row code that is invariant under the shifts of all but
$b$ pairs of a pairing of the entries has at most $q^{d^2/2+b+o(d^2)}$ elements over an alphabet of $q$ symbols, with
an explicit threshold that does not depend on the alphabet or on the relation. So a code of rate above one half must
break the symmetry of a pairing on a positive fraction of the entries.
\end{itemize}

\subsection{How the paper builds on the series}\label{sec:intro-series}

\paragraph{The first paper} supplies the frozen statement, the Tape Coding Conjecture, and the reduction from machines
that are correct on a fraction $2^{-O(d\log d)}$ of the inputs to codes on block networks, which is what makes the
robust form of the conjecture reach machines that are not oblivious~\cite{FirstPaper2026}. \Cref{thm:T} is stated
against that conjecture as it stands there.

\paragraph{The fourth and fourteenth papers} supply the cut. The fourth paper showed that for oblivious machines two
statements together exclude a transposer: every run can be cut cheaply to bounded depth, and index codes on the
resulting systems of cones are long~\cite[Theorem~4.10]{FourthPaper2026}. The fourteenth paper proved the first of the
two for all runs~\cite[Theorem~7.5]{BoundedStack2026}. \Cref{thm:W} starts from that cut, and \cref{thm:T} applies it to a block
network after lifting the network to a walk.

\paragraph{The seventeenth paper,} the synthesis, reduced the oblivious case to one function, the best rate of a
storage code on the cone relations left by cheap cuts, and described the statement that would decide it
(\cref{sec:intro-diagnosis}). It also introduced the codes of this paper, under the name of $\rho$-codes, proved the
bounds $w/(w+1)$ and one half for $w\le2$, and asked for one half at every window~\cite{SeventeenthPaper2026}. The
present paper shows that these codes, and a far weaker bound than one half, are all that the question about machines
needs.

\paragraph{Other papers.} The storage-code bounds of the fifth and seventh papers are used for the classes in which
the conjecture is known and for the certificates of \cref{thm:E}~\cite{DirectedRecovery2026,SeventhPaper2026}. The
thirteenth and fifteenth papers constructed storage codes of rate tending to one on abstract relations of bounded
degree~\cite{StorageCodes2026,HighRate2026}; whether such codes can be whole-row codes is one of the open questions of
\cref{sec:open}. The eighteenth paper treated machines that are not oblivious through their own
runs~\cite{AdaptiveTransfer2026}; \cref{thm:T} reaches them through the first paper instead.

\subsection{The diagnosis of the synthesis, point by point}\label{sec:intro-diagnosis}

The synthesis listed five properties that a statement deciding the oblivious case must
have~\cite[Section~4.4]{SeventeenthPaper2026}. We go through them for \cref{conj:MCC}.
\begin{enumerate}[label=(\arabic*),leftmargin=2em]
\item \emph{It must tie what is known to where it can be, for all schedules.} After \cref{thm:W} the place is a set of
whole rows, and the schedule has disappeared from the statement: the conjecture quantifies over all~$\rho$.
\item \emph{It must fail for abstract separated relations, on which codes of rate tending to one exist.} The
conjecture says nothing about such relations. It concerns whole-row relations, in which an entry is helped by $wd$
others, all of them in $w$ rows. No code of rate above one half is known on these at a fixed window as $d$ grows.
\item \emph{It must fail for motion that ignores the data.} The conjecture is not about motion. The codes of
\cref{thm:W,thm:T} are extracted from correct outputs, and a machine that moves its heads and transposes nothing yields
none.
\item \emph{It must constrain all scales at once.} The conjecture is a statement at one scale, and the scales enter
through the window: a cheaper cut leaves deeper cones and a larger window, so the deficit must not vanish as the
window grows. \Cref{thm:F} makes the exchange rate explicit.
\item \emph{It must use the shape of tapes,} because a bound on coding gains for graphs of bounded degree in general
would be a circuit lower bound. The shape of tapes is used in the reduction: the cheap cut exists because heads walk,
and cones meet few rows because the input is read row by row. What is left in the statement is the shape of the
array: a row and a column meet in one entry.
\end{enumerate}
So the diagnosis is met in a different way than the synthesis expected in points~(4) and~(5), and the restriction to
oblivious machines, under which it was made, is removed. The wall has not fallen. It has been moved to a statement
about arrays, and \cref{thm:O,thm:E,thm:I} say what is known on the other side.

\subsection{Questions of earlier papers}\label{sec:intro-questions}

\begin{itemize}[leftmargin=1.5em]
\item The fourth paper remarked that no implication was known, in either direction, between the plain Tape Coding
Conjecture and the pair of statements of its decomposition~\cite[Remark~4.12]{FourthPaper2026}. That question stays
open. What is new is one statement that implies both: \cref{conj:MCC} gives the bound on index codes at every depth
(\cref{cor:W-obl}) and the Tape Coding Conjecture (\cref{thm:T}).
\item The synthesis asked whether whole-row codes have rate one half at every window, over a fixed
alphabet~\cite[Problem~9.4(c)]{SeventeenthPaper2026}. We do not answer this. \Cref{thm:T,thm:F} say what any bound
below one, at the alphabet of the transposition problem, would give, and \cref{thm:I} proves one half for codes with
the symmetry of a pairing.
\item The synthesis observed that for small windows the acyclic-subgraph bound might still suffice, and that
computations indicated that for large windows it does not~\cite[Section~9]{SeventeenthPaper2026}. \Cref{thm:O} proves
the second half: on mixing relations the best acyclic bound has rate tending to one as the window grows. Window three
is still open.
\end{itemize}

\subsection{What is left}\label{sec:intro-left}

One statement: \cref{conj:weak}, or the stronger \cref{conj:MCC}, or the still stronger \cref{conj:half}. The weak form
lets the deficit decrease as the window grows, but only very slowly (\cref{rem:weak-demand}), while the bounds that
hold for all codes give a deficit of $1/(w+1)$. The smallest open case
of the half-rate conjecture is window three on relations whose rows and columns are spread. No converse is known: a
family of whole-row codes of rate tending to one would refute the conjecture and would not be a machine.
\Cref{sec:construction} says what such a family, and what a machine, would have to look like, and \cref{sec:open}
states the open problems.

\subsection{Formal status, and how to read the paper}\label{sec:intro-formal}

\Cref{thm:W,thm:T,thm:F,thm:N,thm:O,thm:I} and the finite and conditional parts of \cref{thm:E} are formalized in
Lean~4 with Mathlib~\cite{Mathlib2026}; each carries a line naming its declarations, and \cref{sec:leanindex} lists
them. They use the axioms \lean{propext}, \lean{Classical.choice} and \lean{Quot.sound} only. \Cref{thm:T,thm:F}
conclude the frozen proposition of the first paper. Not formalized, and proved in the text: the two existence statements
\cref{prop:mixing,prop:lifts}, the exclusion of \cref{prop:hyperfinite}, and the code that attains the Fano optimum.
Remarks are not formalized either: in particular the arithmetic of \cref{rem:weak-demand} and the sketch of
\cref{rem:markov}; \cref{prop:known} and the results in the table of \cref{sec:invariance} are cited from earlier
papers with their status there. A
reader who wants only the reduction can read \cref{sec:prelim,sec:rows,sec:tape,sec:weak}; a reader who wants to work
on the conjecture can start at \cref{sec:normal}.

\subsection{Related work}\label{sec:intro-related}

\paragraph{Codes.} Whole-row codes are storage codes, or equivalently index codes with side information, on a
particular family of relations; for these notions see~\cite{BBJK2011,Riis2007}. For storage codes on graphs it is
known that the simplest bounds are not always attained and that natural conjectures about optimal codes
fail~\cite{ChristofidesMarkstrom2011,Baber2017}; \cref{thm:E} is a statement of this kind for whole-row relations. In
coding theory a \emph{matrix code} usually means a code in the rank metric and an \emph{array code} a code for
arrays of disks. The conjecture of this paper is named after the matrix that is transposed and concerns neither.

\paragraph{Cuts.} Fixing the values on a cut of a computation of small depth goes back to Valiant's depth
reduction~\cite{Valiant1977}; Paul, Pippenger, Szemer\'edi and Trotter cut the computations of multitape
machines~\cite{PPST1983}. With their technique Mittal and Raz showed that a function with a static property, block
rigidity, cannot be computed in linear time by a multitape machine, even with advice, and that an explicit such
function would follow from a strong parallel repetition theorem~\cite{MittalRaz2021}; a conjecture of theirs about a
game that predicts a transpose is known to fail~\cite{Mittal2025}. Our reduction has the same shape, a static
statement that excludes linear time through a cut, with a different statement: it is about codes on arrays, and it
is specific to transposition.

\paragraph{Network coding.} Adler, Harvey, Jain, Kleinberg and Lehman related transposition in the input--output
model to network coding~\cite{AHJKL2006}, and the network coding conjecture of Li and Li~\cite{LiLi2004} would have
implied lower bounds for sorting in external memory~\cite{FHLS2019}, for multiplication
circuits~\cite{AFKL2019} and for oblivious machines with two tapes~\cite{AHJKL2006}. That conjecture has been
refuted~\cite{BravermanHe2026,ZhangLiLi2026}. The Tape Coding Conjecture of the first paper restricts it to the
networks of tapes; \cref{conj:MCC} is not a statement about networks at all. Lower bounds for transposition are known
on machines with one work tape~\cite{DMS1991,DietzfelbingerMaass1993}.

\paragraph{What we did not find.} We searched for the whole-row relations, under names such as array codes with
recovery of columns from rows and locally recoverable array codes, and for a statement like \cref{thm:T} among the
works above and those that cite them. We found neither. The search was not exhaustive; in particular we have not read
every dissertation in the area.

\subsection{Organization}\label{sec:intro-org}

\Cref{sec:prelim} fixes the model, the codes and the notation. \Cref{sec:rows,sec:tape,sec:weak} prove the reduction.
\Cref{sec:normal,sec:order,sec:entropy,sec:invariance} study the conjecture: a normal form, what order cannot do, what
entropy can, and where one half is proved. \Cref{sec:construction} is the construction side, \cref{sec:open} the open
problems, and the appendices index the Lean declarations and describe the verification.

\section{Preliminaries}\label{sec:prelim}

\subsection{Machines and the statement}\label{sec:prelim-machines}

We use the model and the formal statement of the first paper~\cite[Section~2]{FirstPaper2026}. A machine has a
read-only input tape, $k$ work tapes and a write-only output tape. A $d\times d$ matrix $A$ with entries of
$b=b(d)=\max(1,\lceil\log_2d\rceil)$ bits is encoded row by row as a word of $L=d^2b$ bits, and the machine must append
the encoding of the transpose. Statement~(T) is the Lean proposition \lean{Transpose.LinearTimeTransposable}: some
machine and some constant $c$ do this within $c(L+1)$ steps for all $d\ge1$ and all~$A$. A machine is
\emph{oblivious} at size~$d$ if the positions of all heads and the times at which output bits are appended are the
same for all inputs of that size.

\subsection{Whole-row codes}\label{sec:prelim-codes}

Rows and columns of a $d\times d$ array are indexed by $[d]=\{0,\dots,d-1\}$, and an array over an alphabet~$\Sigma$
is an element of $\Sigma^{[d]\times[d]}$. A \emph{helper map} is a map $\rho$ from $[d]$ to the subsets of $[d]$; its
\emph{window} is $\max_j\card{\rho(j)}$. We read $\rho(j)$ as the set of rows that help column~$j$.

\begin{definition}[Whole-row codes]\label{def:wholerow}
\formal{StatementsRowSpread.lean}{RowSpread.wholeRowHelpers, RowSpread.IsWholeRowCodeExcept, RowSpread.WholeRowCardinalityBound, RowSpread.UniformWholeRowRateCanonicalBits}
A set $C\subseteq\Sigma^{[d]\times[d]}$ is a \emph{whole-row code for~$\rho$} if for every $j\in[d]$ any two arrays of
$C$ that agree in all entries of the rows $\rho(j)$ agree in all entries of column~$j$. It is \emph{two-sided} if in
addition any two arrays of $C$ that agree in all entries of the columns $\rho(j)$ agree in all entries of row~$j$.
\end{definition}

The function that computes a column from its helper rows is arbitrary, and differs from column to column. An entry
$(i,j)$ with $i\in\rho(j)$ lies in one of its own helper rows; we call it \emph{incident}. There are at most $wd$
incident entries, so they do not affect rates. These are the $\rho$-codes of the
synthesis~\cite[Section~6.7]{SeventeenthPaper2026}.

We also read a helper map as the relation $\{(i,j):i\in\rho(j)\}$ on $[d]$, and as the bipartite graph that joins
row~$i$ to column~$j$ when $i\in\rho(j)$. The relation is \emph{symmetric} if $i\in\rho(j)$ exactly when
$j\in\rho(i)$, and \emph{reflexive} if $j\in\rho(j)$ for all~$j$. Its \emph{incidence girth} is the girth of this
bipartite graph. An \emph{empty $r\times r$ incidence rectangle} is a set $I$ of $r$ rows together with a set $J$ of
$r$ columns such that $i\notin\rho(j)$ for all $i\in I$ and $j\in J$.

\begin{definition}[Rates]\label{def:rate}
The \emph{rate} of a nonempty code $C$ over $\Sigma$ is $\log_{\card\Sigma}\card C/d^2$. For $w\ge1$ let $R(w)$ be the
upper limit, as $d\to\infty$, of the largest rate of a whole-row code of window at most $w$ on $d\times d$ arrays with
entries in $\{0,1\}^{b(d)}$, the \emph{canonical alphabet}.
\end{definition}

\begin{example}[Pairing]\label{ex:pairing}
For $\rho(j)=\{j\}$ the symmetric arrays form a two-sided whole-row code of rate $\tfrac12+\tfrac1{2d}$. So
$R(w)\ge\tfrac12$ for every~$w$.
\end{example}

In this notation \cref{conj:MCC} says $\sup_wR(w)<1$; its Lean form is
\lean{Wall.CommonDeficit}. The Lean form bounds the size by $2^{(1-\delta+\eta)L}$ for every $\eta>0$ and all large~$d$;
this is equivalent to the form of \cref{conj:MCC}, with $\delta$ halved in one direction.

\begin{conjecture}[Half-rate conjecture]\label{conj:half}
\formal{StatementsRowSpread.lean}{RowSpread.WholeRowHalfRateCanonicalBits}
$R(w)=\tfrac12$ for every $w\ge1$.
\end{conjecture}

\begin{proposition}[Known bounds]\label{prop:known}
\partformal{KnownBounds.lean}{KnownBounds.cardinality_pow, KnownBounds.cardinality_bound, KnownBounds.window_one, KnownBounds.cardinality_pow_threshold, KnownBounds.uniform_in_alphabet}{The sharper window-two bound remains the written proof of \cite[Proposition~6.22(b)]{SeventeenthPaper2026}.}
For every alphabet, a whole-row code of window $w$ has rate at most $\tfrac{w}{w+1}(1+\tfrac1d)$, and for $w\le2$ it
has at most $\card\Sigma^{d(d+1)/2+d-1}$ elements. In particular $R(1)=R(2)=\tfrac12$ and
$\tfrac12\le R(w)\le\tfrac{w}{w+1}$.
\end{proposition}

\begin{outline}
The averaging argument is that of \cite[Proposition~6.24(d)]{SeventeenthPaper2026}. Fix a row order and retain, in
each column, the entries up to its last helper row. These entries determine the array by induction in that order.
For a set of $k$ helper rows and one other row, each of the $k+1$ rows is last in equally many orders: exchanging
two row labels is a bijection between the two classes of orders. Summing these counts gives an order of cost at most
$wd(d+1)/(w+1)$. The Lean proof retains the sharper numerator $wd^2+\sum_j\card{\rho(j)}$, and applies to rectangular
arrays as well. In particular
$\card C^{w+1}\le\card\Sigma^{wd(d+1)}$, with no linearity or restriction on the recovery functions.
At window one this gives $\card C\le\card\Sigma^{d(d+1)/2}$; pairing attains it. The threshold for the limiting
upper rate is independent of the alphabet. Window two uses the separate multigraph-ranking argument cited above.
\end{outline}

\begin{remark}[Windows that grow with $d$]\label{rem:growing}
\partformal{GraphMatchingCode.lean}{GraphMatchingCode.exists_code, GraphMatchingCode.mem_code_iff, GraphMatchingCode.cardinality, GraphMatchingCode.isWholeRowCode, GraphMatchingCode.window_iff, GraphMatchingCode.update_nonedge, GraphMatchingCode.mem_linearCode_iff_code, GraphMatchingCode.finrank_linearCode_add}{The graph-existence results and their asymptotic parameters, cited below. The finite graph-to-code conversion is formal.}
In \cref{def:rate,conj:MCC} the window is fixed before $d$ grows. This cannot be relaxed much. Mittal and Raz
conjectured that for some constant $c>0$ no set of $2^{(1-c)d^2}$ matrices of bits has every column a function of
$cd$ rows; in their terms, that the transpose is a block-rigid function~\cite[Conjecture~36]{MittalRaz2021}. In the
terms of this paper that is a deficit $c$ at window $cd$, for entries of one bit. It is false. By a construction of
Braverman and Forbes, given in Mittal's dissertation, for every $\eps>0$ there is an $\alpha>0$ such that for
infinitely many~$d$ some helper map of window at most $d^{1-\alpha}$ has an $\F_2$-linear whole-row code of
codimension at most $d^{1+\eps}$~\cite[Question~10.8.2, Proposition~10.8.3 and Lemma~10.8.5]{Mittal2025}. Used for
each bit of the entries, it gives whole-row codes over the canonical alphabet of rate at least $1-d^{\eps-1}$. The
codes are defined by parities over the induced matchings of a dense Ruzsa--Szemer\'edi graph, and the same
constructions are delivery schemes in coded caching~\cite{ShanmugamEtAl2017,YanEtAl2017}; see the second paper of the
series~\cite{SecondPaper2026}. So the best rate at window $d^{1-\alpha}$ tends to one, and \cref{conj:MCC} would be
false for windows of that size. A second construction of Alon, Moitra and Sudakov lowers the window. For
infinitely many~$d$ it is a graph on $d$ vertices, the words of a fixed length with an edge between two words that
agree in few positions, that misses at most $d^{3/2+\eps}$ edges and whose edges are partitioned into
$d^{2-c\eps^3}$ induced matchings~\cite[Claim~2.11, Lemma~2.12, Theorem~2.20 and the paragraph after it]{AMS2013}. All its
vertices miss equally many others, hence at most $2d^{1/2+\eps}$. Take $\rho(j)$ to be $j$ together with the vertices
missed by~$j$; listing each edge of a matching in both orders gives a partition of the non-incident entries as in
\cref{prop:matchings}, and with it whole-row codes of rate at least $1-2d^{-c\eps^3}$ at window at most
$2d^{1/2+\eps}+1$. At a fixed window, codes of this form have rate at most one half in the limit
(\cref{prop:matchings}).
\end{remark}

\begin{outline}[The finite conversion in \cref{rem:growing}]
Let $G$ be a simple graph on $d$ vertices whose edges are partitioned into $m$ nonempty induced matchings.
Set $\rho(j)=\{i:\{i,j\}\notin E(G)\}$, including $j$ itself. Each matching supplies one block of ordered
entries, with both orders of each edge present. Distinct entries in a block are mutual helpers: if their
cross endpoints were adjacent, the induced-subgraph property and the unique neighbour in a matching would
force the two ordered entries to coincide. The blocks therefore form exactly the partition required in
\cref{prop:matchings}.

For any finite abelian group $A$, impose one prescribed sum per block. The resulting whole-row code has
exactly $\card A^{d^2-m}$ words; diagonal and non-edge entries are free. Its helper width at column $j$ is
$d-\deg_G(j)$, so minimum degree at least $d-w$ gives window at most $w$. Over any field the zero-sum code
is a linear subspace of codimension exactly $m$, since the disjoint nonempty checks give a surjective linear
map to the $m$ check values. These are the finite statements used with the cited graph families; no
existence theorem for those families is part of this formalization.
\end{outline}

\subsection{The Tape Coding Conjecture}\label{sec:prelim-tcc}

We recall the objects of the first paper in the form in which \cref{sec:tape} uses
them~\cite[Sections~8 and~9]{FirstPaper2026}. A \emph{time--space diagram} with $k$ work tapes and length~$T$ records,
for every step, the positions of the input head and of the work heads and the length of the output so far; it is
\emph{valid} if every head moves by at most one cell per step. Fix a granularity $G\ge1$ and let $\ell=\lceil db/G\rceil$, the length of a block. The \emph{block network} of
the diagram has as nodes the \emph{slots}, the consecutive stretches of $\ell$ steps, in their order in time; a
\emph{memory edge} from every slot to the next; and a \emph{chord edge} from a slot to a later one whenever some work
head, during the later slot, is in blocks in which it last resided during the earlier one. The rows of the input that
the input head touches during a slot are \emph{attached} to it, and an output bit is produced at the slot in which it
is appended. A memory edge can carry the blocks in which the heads reside and a chord edge one block for every tape
that has the chord; these are the \emph{profile capacities}~\cite[Definitions~8.2 and~8.3]{FirstPaper2026}.
A \emph{causal code} assigns to every
edge a function of the values on the edges into its tail and of the bits of the rows attached there, and to every
output bit a function of the same data at its slot. It \emph{solves complete exchange} on a set of inputs if on those
inputs every edge value fits its capacity and every output bit is the transposed input bit.

\begin{conjecture}[Tape Coding Conjecture, robust form {\cite[Conjecture~9.1]{FirstPaper2026}}]\label{conj:TCC}
For all $k$, $G\ge1$, $c$ and $C_0$ and all sufficiently large~$d$, no causal code solves complete exchange at the
profile capacities on the block network of a valid diagram with $k$ work tapes and length at most $cL$, for a set of
inputs of co-entropy at most $C_0\,d\,b(d)$.
\end{conjecture}

The Lean proposition is \lean{Transpose.Tape.RobustNoCode_TM}. The \emph{plain form}, \lean{NoCode_TM}, asks the same
for codes that are correct on all inputs. The robust form implies that (T)
fails~\cite[Theorem~9.2]{FirstPaper2026}, and the plain form excludes oblivious linear-time
transposers~\cite[Theorem~9.9]{FirstPaper2026}; both implications are Lean-verified there.

\subsection{Notation}\label{sec:prelim-notation}

The letters of the earlier papers collide, so we fix them here.
\begin{center}
\begin{tabular}{@{}ll@{}}
\toprule
$d$, $b=b(d)$, $L=d^2b$ & side length, bits of an entry, bits of an array\\
$\rho$, $w$ & helper map and its window (in the fourth paper $\rho$ is a rate)\\
$U$, $R$ & normalized window; reverse-degree cutoff\\
$\delta$ & deficit: a bound $2^{(1-\delta)L}$ on the size of a code\\
$k$, $c$, $G$ & work tapes, time constant, blocks per row\\
$\ell=\lceil db/G\rceil$ & cells of a block (the $r$ of the first paper)\\
$\theta$, $D$ & density of a cut among the nodes; depth left by the cut\\
$W$ & number of rows that a cone can meet after the cut\\
\bottomrule
\end{tabular}
\end{center}
Logarithms are to base two unless a base is written. A statement holds \emph{for large $d$} if it holds for all $d$
above a threshold that may depend on every other parameter in scope.

\section{From runs to whole rows}\label{sec:rows}

This section proves the reduction for oblivious machines. It is the direct route, and it shows where the rows come
from: the input head reads the matrix row by row, so a short stretch of computation sees few rows.

\subsection{Cuts, cones and index codes}\label{sec:rows-cones}

Fix an oblivious run of length $T\le c(L+1)$ at size~$d$. Its \emph{dependency graph} has the steps as vertices; a
step depends on the preceding step and, for every work tape, on the last earlier step at which that head was on the
same cell~\cite[Section~4]{FourthPaper2026}. A \emph{cut} is a set $\Gamma$ of steps whose values are revealed; its
\emph{depth} is the largest number of vertices on a path that avoids it. The \emph{cone} of an output bit is the set
of input positions read at the steps from which its appending step can be reached without passing through~$\Gamma$.
The revealed values form an index code: two inputs with the same revealed values that agree on the cone of an output
bit agree on that bit of the transpose~\cite[Corollary~4.5]{FourthPaper2026}. Every run of linear length has a cut of
at most $\eps L$ revealed bits and depth bounded in terms of $\eps$, $k$, $c$ and the width of a
step~\cite[Theorem~7.5]{BoundedStack2026}.

\subsection{Theorem W}\label{sec:rows-W}

\begin{theorem}[Theorem W: whole rows]\label{thm:W}
\formal{WindowSpread.lean}{WindowSpread.native_row_spread_prefix, WindowSpread.nativeRowSpreadQuantitative, WindowSpread.nativeRowSpreadExtension}
Let $k$, $c$ and the width of a step be fixed, and let $\eps>0$. There are $w$ and $D$, depending on these and not
on~$d$, such that for all large $d$ every oblivious run of length at most $c(L+1)$ at size $d$ has a cut $\Gamma$ of
depth at most~$D$ that reveals at most $\eps L$ bits, a set $E$ of at most $\eps d$ exceptional rows of the transpose,
and a helper map $\rho$ of window at most~$w$, such that the cone of every output bit in a row $j\notin E$ of the
transpose lies in the rows $\rho(j)$ of the input.
\end{theorem}
\draftnote{Statement to be checked word for word against the Lean proposition; outline of the formal proof (the
counting of windows along the input walk, the guards, the row cover) to be written.}

\begin{corollary}[The oblivious case]\label{cor:W-obl}
\formal{RowSpreadCrux.lean}{RowSpread.native_indexCode_bound_of_wholeRow, RowSpread.uniformWholeRowRateCanonicalBits_crux, RowSpread.uniformWholeRowRateCanonicalBits_excludes, RowSpread.wholeRow_window_bound_excludes}
If \cref{conj:MCC} holds with deficit~$\delta$, then every index code on the cone system of a cut of bounded depth of
an oblivious linear-time run has at least $(\delta-o(1))L$ bits, at every depth, and no oblivious machine transposes
in linear time. A bound at one explicit window, with a deficit above a fixed multiple of the cost of the cut, already
excludes the machines with the corresponding resources.
\end{corollary}
\draftnote{The constants of this statement (the factor in front of $\delta$, the multiple of the cost of the cut in
the one-window form) are to be taken from the Lean declarations. Outline: fix the revealed values and the exceptional
rows; the fibre is a whole-row code for $\rho$ (\lean{entryWord_wholeRow_of_cones}); sum over fibres.}

The synthesis reduced the oblivious case to one function, the best rate of a storage code on the cone relations left
by cheap cuts~\cite[Section~4.1]{SeventeenthPaper2026}. The corollary replaces that function, which depends on the
schedule, by the rate of whole-row codes: if this rate stays below one, no oblivious machine transposes in linear
time.

\section{From block codes to whole rows}\label{sec:tape}

The Tape Coding Conjecture is about codes on block networks, in which a node may compute anything from the rows
attached to it. We cut such a code as we cut a run. A block network is not the dependency graph of a walk, so the cut
of the fourteenth paper does not apply to it directly; the first step is to realize it inside one.

\begin{theorem}[Theorem T: the Matrix Coding Conjecture implies the Tape Coding Conjecture]\label{thm:T}
\formal{WholeRowTCC.lean}{BlockTCC.robustNoCode_of_commonDeficit, BlockTCC.noCode_of_commonDeficit, BlockTCC.not_linearTimeTransposable_of_commonDeficit, BlockTCC.robustNoCode_of_normalizedTwoSidedDeficit}
If \cref{conj:MCC} holds, then \cref{conj:TCC} holds, in its robust and hence in its plain form. Consequently
\cref{conj:MCC} implies that \textup{(T)} fails: \lean{Transpose.LinearTimeTransposable} is false. The same follows
from the conjecture for two-sided codes on symmetric reflexive relations.
\end{theorem}

\begin{outline}
Suppose a causal code solves complete exchange on a set $F$ of inputs of co-entropy at most $C_0db$, on the block
network of a valid diagram of length at most~$cL$.

\emph{Lifting.} Every slot is replaced by a walk of twelve positions that visits the blocks on both sides of a common
pivot, so that every memory edge and every chord edge of the network becomes a path of the dependency graph of a
walk (\leanfile{BlockIntervalWalk.lean}).

\emph{Cutting.} The cut of the fourteenth paper for that walk, projected back to slots, is a set $\Gamma$ of at most
$\theta N$ of the $N$ slots after whose removal every path has bounded length (\leanfile{BlockCut.lean}).

\emph{Rows.} A slot has at most two attached input rows, so the slots from which an output slot can be reached
without passing through $\Gamma$ meet boundedly many rows; rows of the transpose that meet too many output slots are
paid for (\leanfile{BlockRowGrouping.lean}, \leanfile{BlockCodeLocal.lean}).

\emph{Fibres.} Fix the messages on the edges out of $\Gamma$ and the paid rows. On the inputs of $F$ that produce
them, every remaining column of the array is a function of boundedly many rows: a whole-row code
(\leanfile{BlockCodeTransfer.lean}).

\emph{Invoice.} The profile capacities bound the number of fixed bits by a multiple of $\theta L$. Some fibre has at
least $2^{L-\text{cost}}$ elements, and for small $\theta$ this contradicts the deficit
(\leanfile{BlockCutPrice.lean}, \leanfile{BlockCodeInvoice.lean}).
\end{outline}
\draftnote{Give the constants (the price $(4kG+1)\ell$ of a cut slot, the bound on the rows of a cone) and the
treatment of the co-entropy of $F$.}

The theorem is a common strengthening, not a comparison: it does not say that the Tape Coding Conjecture implies
anything about whole-row codes. A whole-row code of rate near one need not be computable by a causal code on any
diagram, and no converse of \cref{thm:T} is known.

\section{The weak form}\label{sec:weak}

\Cref{thm:T} uses one deficit for all windows. The proof uses less: for given resources it needs a bound at one
window, at one large size, with a deficit above the cost of the cut. Stating exactly that gives a weaker conjecture
with the same consequence. By \cref{thm:N} it may be stated for two-sided codes on symmetric reflexive relations.

For $k\ge5$, $G\ge1$, $c\ge4$, $d\ge1$, $\theta>0$ and $R\ge0$ put $b=b(d)$, $L=d^2b$, $\ell=\lceil db/G\rceil$ and
\begin{align*}
  D&=\text{the depth of the cut of density $\theta$ of~\cite[Theorem~7.5]{BoundedStack2026}},&
  W&=2(G+1)(2k+3)^{D},\\
  K&=c+4+5(k-3)(cG+2),&
  M&=(4kG+1)\,\ell\,\bigl\lfloor\theta\lceil cL/\ell\rceil\bigr\rfloor,\\
  h&=\begin{cases}0,&R\ge d,\\ \min\bigl(d,\lfloor Wd/(R+1)\rfloor\bigr),&R<d,\end{cases}&
  X&=2dh-h^2 .
\end{align*}

\begin{conjecture}[Matrix Coding Conjecture, weak form]\label{conj:weak}
\formal{Wall.lean}{Wall.Conjecture}
For all $k\ge5$, $G\ge1$, $c\ge4$ and every $d_0$ there are $d\ge\max(d_0,1)$, $R$, $\theta>0$ and $\delta$ such that
\begin{enumerate}[label=(\roman*),leftmargin=2em]
\item every two-sided whole-row code on $d\times d$ arrays with entries in $\{0,1\}^{b}$, for a symmetric reflexive
helper map of window at most $W+R+1$, has at most $2^{(1-\delta)L}$ elements, and
\item $\delta L>2(Kd+M)+bX$.
\end{enumerate}
\end{conjecture}

\begin{theorem}[Theorem F: the weak form suffices]\label{thm:F}
\formal{Wall.lean}{Wall.of_commonDeficit, Wall.excludes_adaptive, Wall.of_previous, Wall.commonDeficit_implies_previous}
\Cref{conj:MCC} implies \cref{conj:weak}, and \cref{conj:weak} implies that \textup{(T)} fails.
\end{theorem}
\draftnote{Outline: final-output diagrams of co-entropy $Kd$ from the first paper (\lean{codes_coentLinear_at}), the
end invoice with integer budgets (\leanfile{FiniteWholeRowPayments.lean}), the normalization.}

\begin{remark}[What the weak form asks]\label{rem:weak-demand}
Condition~(ii) reads, up to lower-order terms, $\delta>2(4kG+1)c\,\theta+2W/(R+1)$ at window $U=W+R+1$. So a bound
$1-f(U)$ on the rate at window~$U$ need not have $f$ bounded below: it suffices that for some $\theta$ and $R$ the
value $f(W+R+1)$ exceeds the right-hand side. \emph{Heuristic:} the depth $D$ of the cut grows very fast as $\theta$
falls and $W$ is exponential in~$D$, so that $\theta$ is of the order of an inverse power of $\log\log W$; a deficit
that falls more slowly than every such power would then suffice. We have not proved an estimate of this form. On the other hand the general
bound $f(U)=1/(U+1)$ of \cref{prop:known} never meets it, since $1/(W+R+2)<2W/(R+1)$. The symmetric arrays show
$\delta\le(d-1)/(2d)$, and a window of $d$ or more admits all arrays, so every witness has $W+R+1<d$.
\end{remark}

\section{Normal forms}\label{sec:normal}

\begin{theorem}[Theorem N: two-sided codes on symmetric relations]\label{thm:N}
\formal{TwoSidedCodeTransfer.lean}{RowSpread.difference_fibre_twoSided, RowSpread.wholeRow_card_sq_le_of_twoSided, RowSpread.commonDeficit_iff_twoSided}
\formal{NormalizedRateTransfer.lean}{RowSpread.twoSided_card_le_of_fixedCross, RowSpread.commonDeficit_iff_normalizedTwoSided, RowSpread.commonDeficitCanonicalBits_iff_normalizedTwoSided}
Let $C$ be a whole-row code for $\rho$ over an alphabet $\Sigma$ with a group structure.
\begin{enumerate}[label=(\alph*),leftmargin=2em]
\item For some array $z$, the set $C\cap(z+C^{\mathsf T})$ is a two-sided whole-row code for $\rho$ with at least
$\card C^2/\card\Sigma^{d^2}$ elements.
\item If every two-sided code on a symmetric reflexive relation of window $w+R+1$ has at most $\card\Sigma^{a}$
elements, then every two-sided code for a helper map of window $w$ has at most $\card\Sigma^{a+2wd^2/(R+1)}$ elements.
\item \Cref{conj:MCC} is equivalent to the same statement for two-sided codes, and to the same statement for
two-sided codes on symmetric reflexive relations.
\end{enumerate}
\end{theorem}
\draftnote{Outline: every pair of elements of $C$ occurs in exactly one $z$-fibre; freezing the rows and columns of
high reverse degree. Examples: the cyclic relation, where the one-sided excess $d-1$ over pairing becomes~$1$.}

So the conjecture loses nothing if rows are required to be functions of columns as well. A deficit $\delta$ for
two-sided codes gives $\delta/2$ for all codes.

\section{Order and closure}\label{sec:order}

The bounds of \cref{prop:known} come from one argument: a set $A$ of entries is \emph{acyclic} if no entry of $A$ is
reached from itself along helpers inside~$A$, an entry that is its own helper counting as a cycle; then a code has at
most $\card\Sigma^{d^2-\card A}$ elements. We determine the best such bound and show what it cannot do.

\begin{theorem}[Theorem O: order bounds, exactly]\label{thm:O}
\formal{AcyclicOrder.lean}{WholeRow.exists_rowSet_superset, WholeRow.rowSet_card_add_cost, WholeRow.exact_acyclic_card}
\formal{ClosureBarriers.lean}{ClosureBarriers.assigned_rule_card_le, ClosureBarriers.transported_invoice_ge}
\formal{RectangleBarrier.lean}{RectangleBarrier.rectangle_barrier, RectangleBarrier.rectangle_invoice_ge}
Let $\rho$ be a helper map on $[d]$, for arrays with rows $[d]$ and columns in any finite set.
\begin{enumerate}[label=(\alph*),leftmargin=2em]
\item For an order of the rows, the entries that lie strictly after the last helper row of their column form an
acyclic set, and every acyclic set is contained in one of these. The largest acyclic set has $d^2-\sigma(\rho)$
entries, where $\sigma(\rho)$ is the least value, over the orders of the rows, of the sum over the columns of the
position of the last helper row.
\item If the recovery rules of $m$ helper relations together generate all entries from a set of $s$ entries, then
$d^2-s$ is at most the sum of the sizes of their largest acyclic sets. In particular the rules of transported copies
of one relation never give a better bound than its own acyclic bound.
\item If no $r$ rows and $r$ columns are free of incidences of $\rho$, then $\sigma(\rho)\ge(d-r+1)^2$.
\end{enumerate}
\end{theorem}
\draftnote{The inequality of (c) and the form of (b) are to be checked against the Lean declarations. Outlines; the
values $\sigma=41$ for the Fano and grid relations.}

\begin{proposition}[Mixing relations of bounded window]\label{prop:mixing}
\notformalized
For every $\eta>0$ there is a $t$ such that for all large $d$ there is a symmetric reflexive helper map on $[d]$ of
window at most $2t+1$ for which no $\lceil\eta d\rceil$ rows and $\lceil\eta d\rceil$ columns are free of incidences.
\end{proposition}
\draftnote{Complete proof: the union of $t$ random perfect matchings with their inverses and the diagonal, for
$t\eta^2>2\ln2$, and a union bound over pairs of sets.}

By (c) and \cref{prop:mixing} the best acyclic bound has rate at least $(1-\eta)^2$ on some relations of bounded
window, so no bound of this
kind, and by (b) no closure argument with transported copies, proves \cref{conj:MCC}. This confirms the indication of
the synthesis for large windows~\cite[Section~9]{SeventeenthPaper2026}.

\section{Entropy beyond order}\label{sec:entropy}

A \emph{valid polymatroid} for $\rho$ is a normalized, monotone, submodular function $h$ on sets of entries with
$h(\{x\})\le1$ and $h(S\cup\{x\})=h(S)$ whenever $S$ contains the helper rows of the column of~$x$. The entropy of a
uniformly random element of a code, in units of $\log\card\Sigma$, is a valid polymatroid whose value on the set of
all entries is $\log_{\card\Sigma}\card C$. So the largest value of a valid polymatroid on that set, the \emph{Shannon
value} $\mathrm{Sh}(\rho)$, bounds $\log_{\card\Sigma}\card C$~\cite{SeventhPaper2026}.

\begin{theorem}[Theorem E: entropy beats order]\label{thm:E}
\formal{FanoFull.lean}{FanoFull.full_upper, FanoFull.wholeRow_card_cube}
\formal{FanoAttainment.lean}{FanoAttainment.full_attainment, FanoAttainment.full_shannon_exact}
\formal{WholeRowAcyclic.lean}{WholeRowAcyclic.fano_exact_acyclic_card, WholeRowAcyclic.fano_gap, WholeRowAcyclic.grid_exact_acyclic_card, WholeRowAcyclic.grid_gap}
\formal{FanoLift.lean}{FanoLift.conditional_gap, FanoLift.conditional_acyclic_gap, FanoLift.conditional_gap_one_percent, FanoLift.completeLift_gap}
\begin{enumerate}[label=(\alph*),leftmargin=2em]
\item \emph{Fano.} Let $d=7$ and let $\rho(j)$ be the three points of the line $j$ of the Fano plane. Then
$\mathrm{Sh}(\rho)=119/3$, a code over an alphabet of $q^3$ symbols attains it, and $\sigma(\rho)=41$: every
acyclic bound is at least $4/3$ entries above the truth.
\item \emph{Grid.} For the $9\times6$ array whose columns are the six lines of the $3\times3$ grid and whose rows are
its nine points, $\mathrm{Sh}\le204/5$ and $\sigma=41$. The entries outside the incidences have no mutual triangles.
\item \emph{Lifts.} Let $d=7n$, let every row and column of the Fano relation be replaced by $n$ copies, and let
$\rho$ be a helper map whose helper rows are copies of points on the line of the column. If every $r$ copies of a
point and every $r$ copies of a line through it contain a helper incidence, then for every valid polymatroid $h$ and
every acyclic set~$A$,
\[
  d^2-\card A-h(\text{all entries})\ \ge\ 41\,(n-r+1)^2-\tfrac{119}{3}\,n^2 ,
\]
which exceeds $d^2/100$ when $99n\le100(n-r+1)$.
\end{enumerate}
\end{theorem}
\draftnote{Certificates and their generator; the three-lane code; the projection lemma
(\lean{MacroProjection.upper_of_projection}); the threshold argument for every order of the fine rows.}

\begin{proposition}[Mixing lifts exist]\label{prop:lifts}
\notformalized
For every large $n$ there is a lift as in \cref{thm:E}(c) with $r=\lceil n/100\rceil$, with at most $3600$ helper
rows for every column and at most $3600$ columns helped by every row, and with incidence girth at least
$c\log n$ for a constant $c>0$.
\end{proposition}
\draftnote{Complete proof: unions of $1200$ random perfect matchings on every incidence; the union bound for empty
rectangles; deletion of one vertex on every short cycle.}

These relations have window $3600$, and $\sigma\le(41/49)d^2$ already bounds their codes away from rate one. The
theorem is evidence about methods, not about the conjecture: inequalities that are not of order type can hold at
positive density at a fixed window.

\section{Where one half is proved}\label{sec:invariance}

One half is known for $w\le2$ (\cref{prop:known}), for relations whose rows or offsets are not
spread~\cite{SeventeenthPaper2026}, and for several classes of relations and codes from the earlier papers, listed in
\cref{tab:classes}. The theorem of this section adds a class defined by a symmetry of the code, with no condition on
the relation.

Let $\iota$ be an involution of the entries, for example $(i,j)\mapsto(j,i)$. A \emph{pairing shift} adds one group
element to the two entries of an orbit of~$\iota$, or to a fixed entry.

\begin{theorem}[Theorem I: invariance under a pairing]\label{thm:I}
\formal{WholeRow/Pairing.lean}{WholeRow.pairing_card_sq_le_forward, WholeRow.pairing_card_pow_le_half}
\formal{WholeRow/ParityTrimming.lean}{WholeRow.parity_card_le_forward, WholeRow.parity_card_pow_le_baseline}
Let $\Sigma$ be a finite abelian group, $\rho$ a helper map of window $w$ on $[d]$, and $C$ a whole-row code
for~$\rho$ that is invariant under the pairing shifts of all orbits of $\iota$ except $b$ orbits of two entries. Then
for every $n\ge1$ with $d\ge n(w+2w^3n)$,
\[
  \card C^{2n}\ \le\ \card\Sigma^{\,n(d^2+2b)+2d^2},
\]
that is, the rate is at most $\tfrac12+b/d^2+1/n$. If in addition the shifts of $u$ fixed entries that are not
incident are missing, the bound holds with $2b+u$ in place of~$2b$. For codes invariant under the shifts of $m$
disjoint parity blocks of at least two entries each, the rate is at most $1-m/d^2+1/n$.
\end{theorem}
\draftnote{The statement, in particular the two variants, is to be checked against the Lean declarations. Outline:
compatibility from recovery; records separate cosets; the count; deletion of heavy rows at one
symbol per block; lanes. Table of the classes with one half, with citations. These declarations are in the shared
library; the verification appendix must say which build covers them.}

The theorem excludes one mechanism at every window: a code above the pairing rate cannot keep the pairing symmetry on
all but $o(d^2)$ orbits. Nothing is known to force that symmetry on the codes that come from machines.

\section{The construction side}\label{sec:construction}

What would a counterexample to the conjecture, and what would a linear-time transposer, have to look like?

\subsection{What a code of rate near one must avoid}\label{sec:constr-code}

Suppose that whole-row codes of rate $1-\eps$ exist at a window that depends on $\eps$ only, for all large~$d$.
\begin{itemize}[leftmargin=1.5em]
\item By \cref{thm:N} there are such codes, of rate $1-2\eps$, that are two-sided on symmetric reflexive relations.
\item By \cref{prop:known} and the classes of \cref{sec:invariance} the window is at least three and the rows and
offsets of the relation are spread.
\item By \cref{thm:I}, for every involution of the entries the codes fail to be invariant under many of its pairing
shifts: if the shifts of all fixed entries are kept, rate $1-\eps$ forces the number $b$ of orbits of two entries
without invariance to be at least $(\tfrac12-\eps-o(1))d^2$, and an involution has at most $d^2/2$ such orbits, so
invariance survives on a fraction of at most about $2\eps$ of them. Fixed entries that lose their shifts are charged
separately, by the term $u$ of the theorem.
\item The examples above the pairing rate that are known at small sizes, such as the Fano code of \cref{thm:E}, gain
from the incident entries and from short cycles of helpers. At a fixed window the incident entries are a fraction
$O(w/d)$ of all entries, so a gain of positive density as $d$ grows cannot come from them alone.
\item \Cref{thm:O} gives no information on mixing relations. \Cref{thm:E} shows that this is a limit of the method
and not evidence for such codes: entropy inequalities can hold where acyclic bounds do not.
\end{itemize}
The thirteenth paper constructs storage codes above half rate on abstract relations of bounded degree, and the
fifteenth amplifies them to rate tending to one~\cite{StorageCodes2026,HighRate2026}. In a whole-row relation an entry has $wd$ helpers, which fill $w$ rows,
and we do not know whether those codes can be laid out in this way (\cref{prob:highrate}).
Whole-row codes of rate tending to one do exist when the window may grow like $d^{1-\alpha}$ (\cref{rem:growing}).
They are not the codes asked for above, whose window depends on $\eps$ only. They show that a proof of the conjecture
cannot be indifferent to the size of the window (\cref{prob:growing}).

The codes of \cref{rem:growing} have one form, and at a fixed window that form gives no more than the symmetric
arrays. Call a set $M$ of non-incident entries an \emph{induced matching} for~$\rho$ if any two of its entries
$(i,j)\ne(i',j')$ satisfy $i\in\rho(j')$ and $i'\in\rho(j)$. Its entries then lie in distinct rows and in distinct
columns. Let $\mathcal M$ be a partition of the set of non-incident entries into induced matchings and let $A$ be a
finite abelian group. Let $C(\mathcal M)$ be the set of arrays over $A$ whose entries sum to zero over every
$M\in\mathcal M$. It is a whole-row code for~$\rho$: an entry $(i,j)$ of $M$ is minus the sum of the other entries
of~$M$, and these lie in rows of~$\rho(j)$. For $\rho(j)=\{j\}$ and the matchings $\{(i,j),(j,i)\}$ this is, over
$\F_2$, the code of symmetric arrays. The codes of Braverman and Forbes are of this form over $\F_2$, with the induced
matchings of a Ruzsa--Szemer\'edi graph, and so are the delivery schemes of a placement delivery array with as many
packets as users~\cite{YanEtAl2017}.

\begin{proposition}[Parities over induced matchings]\label{prop:matchings}
\notformalized
Let $\rho$ have window $w$ on $[d]$ and let $\mathcal M$ be a partition of its non-incident entries into induced
matchings. Then at most $(w^2+\tfrac32w)\,d^{3/2}$ entries lie in matchings with three or more entries, and
\[
 \card{C(\mathcal M)}\le\card A^{\,d^2/2+wd/2+(w^2/2+3w/4)\,d^{3/2}} .
\]
So at a fixed window the codes of this form have rate at most $\tfrac12+O(d^{-1/2})$.
\end{proposition}
\begin{informalproof}
Let $\ell(i)$ be the number of columns $j$ with $i\in\rho(j)$, so that $\sum_i\ell(i)\le wd$. Call a row \emph{heavy}
if $\ell(i)>\sqrt d$. There are fewer than $w\sqrt d$ heavy rows, and they contain fewer than $H=wd^{3/2}$ entries.

Let $e_1=(a_1,b_1)$ be an entry of a matching $M\in\mathcal M$ with at least three entries. Suppose first that $M$
has an entry $e_2=(a_2,b_2)\ne e_1$ in a row that is not heavy, and let $e_3=(a_3,b_3)$ be a third entry of~$M$. Then
$a_2\in\rho(b_1)$, $a_2\in\rho(b_3)$ and $a_1\in\rho(b_3)$. The quadruple $(b_1,a_2,b_3,a_1)$ determines $e_1$; hence
$M$, because $\mathcal M$ is a partition; hence $e_2$, the entry of $M$ in row~$a_2$, and $e_3$, the entry of $M$ in
column~$b_3$. There are at most $d$ choices of $b_1$, then $w$ of $a_2\in\rho(b_1)$, then $\ell(a_2)\le\sqrt d$ of
$b_3$, then $w$ of $a_1\in\rho(b_3)$. So at most $w^2d^{3/2}$ entries $e_1$ are of this kind. In the other case all
entries of $M$ other than $e_1$ lie in heavy rows. Then either $e_1$ lies in a heavy row, or $e_1$ is the only entry
of $M$ outside the heavy rows and $M$ has at least two entries in heavy rows. There are fewer than $H$ entries of the
first kind and fewer than $H/2$ of the second. So at most $B=(w^2+\tfrac32w)d^{3/2}$ entries lie in matchings with
three or more entries.

There are at least $d^2-wd$ non-incident entries, of which at least $d^2-wd-B$ lie in matchings with at most two
entries. So $\card{\mathcal M}\ge(d^2-wd-B)/2$. The matchings are disjoint and each imposes one equation, so
$\card{C(\mathcal M)}=\card A^{\,d^2-\card{\mathcal M}}$.
\end{informalproof}

The proof uses that the window is fixed: the bound is vacuous once $w$ is of order $d^{1/4}$. For matchings that all
have the same size $r\ge3$, Alon, Moitra and Sudakov state a stronger bound, which in the present terms reads
$wd\ge\Omega(r^{2/3}d^{4/3})$~\cite[Theorem~3.9]{AMS2013}; its proof is sketched there and also follows paths of
length three among the incident pairs. \Cref{prop:matchings} says nothing about linear codes whose equations are not
of this mutual form, nor about nonlinear codes. For those the question is the conjecture.

\subsection{What the conjecture forgets about machines}\label{sec:constr-forgets}

Between a machine and \cref{conj:MCC} lie several statements, each obtained from the one before by forgetting
something: the run as a table of steps, which is a restatement of the machine
(\lean{SameTable.uniformTableaux_iff}); the values on a cut with the few records that each output can reach; an index
code for all inputs on the cone system; the whole-row codes on its fibres. The reduction of \cref{sec:rows} never uses
that all steps apply the same rule.

\begin{proposition}[Runs that a cheap cut splits]\label{prop:hyperfinite}
\notformalized
Let an oblivious machine with $k\ge1$ work tapes run in time $c(L+1)$ with steps of $s$ bits. Suppose that for some
$\theta$ with $s(k+2)c\,\theta<\tfrac12$ and all large $d$ there is a set of at most $\theta T$ steps whose removal
leaves components of the dependency graph of \cref{sec:rows-cones}, edges taken without direction, with
$o(\log\log d)$ output steps each.
Then the machine does not transpose.
\end{proposition}
\begin{informalproof}
The step width $s$ bounds the joint record at a step, including the state and the work symbols needed to continue
the run. Write $n(d)\ge1$ for an integer upper bound on the outputs per component; replacing it by a larger integer
does not change $n(d)=o(\log\log d)$. Choose constants $0<\delta<1$ and $\eta>0$ such that
\[
 (k+2)(s\theta c+\eta)<(1-\delta)/2.
\]
In addition to the assumed cut $\Gamma_h$, use the depth reduction for oblivious machines~\cite[Corollary~9.3]{BoundedStack2026} to obtain a cut
$\Gamma_D$ of record cost at most $\eta L$ and a remaining depth $D$ independent of $d$. For
$\Gamma=\Gamma_h\cup\Gamma_D$ the undirected components only shrink, both conclusions hold, and
\begin{equation}\label{eq:fragment-upper}
 s\card\Gamma\le s\theta c(L+1)+\eta L.
\end{equation}
Because the schedule is oblivious these cuts and the ensuing helper sets can be fixed before the input.

Each output now has at most $B=\sum_{j=0}^{D}(k+1)^j$ raw input helpers. Apply the paid degree reduction of
\cite[Corollary~3.8]{DirectedRecovery2026}, conditioning on input coordinates that occur in too many cones.
For a fixed sufficiently large reverse-degree cutoff this retains at least $(1-\delta)L$ demanded coordinates,
and their paired-port incidence relation has both degrees at most a constant $\Delta\ge2$. The native
paired-distance estimate of \cite[Theorem~3.3]{DirectedRecovery2026} is needed with a slowly growing radius here,
not merely in its fixed-radius asymptotic form. Its finite packing proof gives, at radius $2n-1$,
\[
 \#\{\text{retained targets at paired distance at most }2n-1\}
 \ \le\ V_{2n-1}\!\sum_{j=1}^{n}2j\,\bigl(a(2j-1)+3\bigr)
              \frac{L}{\log_2d},
\]
where $V_r=\sum_{j=0}^{r}\Delta^j$ and $a=4c_*(B+1)(k+2)(k+6)$. Here $c_*=2c$ ensures
$T\le c_*L$ for $L\ge1$, as required by that finite estimate. The coefficient is
$O_{k,D,c,\Delta}(n^3\Delta^{2n})$. Since $n=o(\log\log d)$, its ratio to $\log_2d$ tends to zero. Remove these
$o(L)$ targets, leaving at least $(1-\delta)L-o(L)$ good ones. This step uses the explicit transport estimate of the fifth paper~\cite{DirectedRecovery2026}, whose constants are
proved there in the text and are not formalized; only its finite packing step is in Lean.

Partition the good outputs according to the uncut component containing their output step, with a singleton patch
for a good output at a cut root. A patch has at most $n$ targets. Its induced paired-port graph has at most $2n$
vertices. If the two ports of a target were connected there, a shortest path would have length at most $2n-1$,
contrary to its being good. Thus the separated-components lemma~\cite{DirectedRecovery2026} bounds every storage
code on the patch's induced relation by $2^{|J|/2}$.

For each patch fix all exterior input bits, including the paid conditioning coordinates. Let $T_J$ be the cut
records at its incoming boundary; include the root record for a cut-root output. Local determination on the DAG
makes these records an index code for the demands in $J$, with the induced raw helpers. The fibre bound gives
$s|T_J|\ge |J|/2$. The records need not be independent and are not charged as a joint entropy. This is precisely
the abstract local-view invoice, formalized as \lean{LocalRecordViews.half_invoice}.

A cut vertex has at most $k+1$ forward children, each entering at most one uncut component, and at most one output
of its own. Hence $\sum_J|T_J|\le(k+2)|\Gamma|$. Summing the patch inequalities gives
\[
 s(k+2)|\Gamma|\ge(1-\delta)L/2-o(L),
\]
contradicting \eqref{eq:fragment-upper} for large $d$. The argument uses outputs per component, not its number of
steps or its time span. No implication in either direction with the macroscopic-motion classes of
\cite{FourthPaper2026} is established.
\end{informalproof}


\begin{remark}[Fibres are nearly Markov]\label{rem:markov}
Proof sketch only: this remark is not used elsewhere and has not been independently reviewed. In the setting of
\cref{thm:T}, fix a value $z$ of the messages on the edges out of the cut. Let $G_z$ be the set of all inputs for which
every cut slot, computed from the prescribed incoming messages and from its attached rows, emits the prescribed
outgoing messages. Each of these conditions involves boundedly many rows, so the uniform distribution on $G_z$
factorizes over the rows and has the Markov property of the hypergraph of these conditions. The code of
\cref{thm:T} is $C_z=F\cap G_z$, where $F$ is the set of inputs on which the block code is correct, of co-entropy~$a$.
A counting argument gives a $z$ for which $C_z$ is large and
$D\bigl(\mathrm{Unif}(C_z)\,\big\|\,\mathrm{Unif}(G_z)\bigr)\le a+2$ bits. For such a $z$, if a set $S$ of rows separates
two sets of rows in that hypergraph, their mutual information given the rows $S$ is at most $a+2$ on $C_z$, and
anything that is a function of each of the two sets has conditional entropy at most $a+2$ given~$S$. An entry is a
function of its own row and of the helper rows of its column, so this bounds what helper rows can know across a
separator. The bound is paid once for every separator, and we do not know how to combine separators.
\end{remark}

\subsection{What a machine would need}\label{sec:constr-machine}

By the proof of \cref{thm:T}, read in the other direction, a linear-time transposer yields for every $\eps>0$ whole-row codes of
rate at least $1-\eps$ at a window that depends on $\eps$ and on the resources of the machine only, for all large~$d$;
for an oblivious machine \cref{thm:W} produces them directly from a cut of its run. So a construction of a machine
contains a refutation of \cref{conj:MCC}, with everything listed in \cref{sec:constr-code}. The converse fails as far
as we know: such codes are sets, and a machine must also produce them by passes over tapes in linear time, with head
motion that meets the conditions of the earlier papers~\cite[Corollary~3.24]{FourthPaper2026}. This is why a
refutation of the conjecture would close this route to an impossibility proof without deciding~(T).

\section{What is left}\label{sec:open}

\begin{problem}[The conjecture]\label{prob:conj}
Prove or refute \cref{conj:weak}. A proof decides \textup{(T)} in the negative. A refutation is a family of
two-sided whole-row codes whose deficit at window $U$ falls below the bound of \cref{rem:weak-demand} for all large
sizes; it removes this route and does not give a machine.
\end{problem}

\begin{problem}[Window three]\label{prob:three}
Is $R(3)=\tfrac12$? The known bounds are $\tfrac12\le R(3)\le\tfrac34$. The open relations have spread rows and
offsets; on those whose incidence graph has girth above eight no code above the pairing rate and no bound below the
acyclic one is known.
\end{problem}

\begin{problem}[A deficit that does not depend on the window]\label{prob:common}
Find any $\delta>0$ and any class of relations, containing relations with $\sigma(\rho)\ge(1-\eta)d^2$ for every
$\eta>0$, on which every valid polymatroid has value at most $(1-\delta)d^2$.
\end{problem}

\begin{problem}[High-rate codes]\label{prob:highrate}
Can the storage codes of rate tending to one of~\cite{StorageCodes2026,HighRate2026} be whole-row codes at a window
that depends on the rate only?
\end{problem}

\begin{problem}[Converse]\label{prob:converse}
Does the Tape Coding Conjecture imply any bound on whole-row codes? Does a family of whole-row codes of rate tending
to one, with any further property, yield a causal code on a tape diagram?
\end{problem}
\draftnote{Add: the questions of the earlier papers that remain open as they were (the remark of the fourth paper;
Problem~9.4 of the synthesis).}


\appendix
\section{Index of Lean declarations}\label{sec:leanindex}
\draftnote{To be generated from the sources by the index tool once the library of the paper is packaged.}

\section{Building and verifying}\label{sec:verification}
\draftnote{Toolchain and pins; build and axiom commands; which declarations are in the shared library; the frozen
files and their hashes; the literature search; heartbeat settings.}


\bibliographystyle{plainurl}
\bibliography{references}

\end{document}
